\section*{Data}

I will use a quarterly firm panel from Q1 2022 to Q4 2025. I use quarterly data because Helene happened at the end of September 2024, so yearly data would mix months before and after the storm. The period gives about 11 quarters before Helene and 5 quarters after. This is needed to check that treated and control firms moved in a similar way before the shock, and to see if the effect after lasts or fades. I start in 2022 to avoid the COVID years, when R\&D and cash were very unusual. I stop at the end of 2025 to have enough time after but not pick up much later shocks.

Firm accounting data come from Compustat Quarterly, accessed through WRDS. Compustat collects the quarterly reports of all US listed firms, so it gives me R\&D, cash, assets and the other variables for every quarter. The main variable is R\&D intensity, measured as xrdq divided by atq. Xrdq is R\&D spending in the quarter and atq is total assets in the quarter. I scale by assets and not by sales because sales itself can fall after a hurricane when shops close or supply chains break. That would move the ratio even if R\&D did not change. I use the same denominator as for cash, which is cheq divided by atq, so the two results are easy to compare. When xrdq is missing I set it to zero, because Compustat leaves it empty when firms have no real R\&D to report. I also control separately for firms that never report R\&D. I drop financials (SIC 6000--6999) and utilities (SIC 4900--4999) and winsorize all ratios at 1\% and 99\%, because a few very small firms can have extreme values.

Treatment comes from FEMA Disaster Declarations for Helene and the NHC report on the storm path. FEMA is the official list that says which counties got federal help after Helene. A firm is treated if its headquarters county got such a declaration. Headquarters county comes from Compustat header data (addzip) and is checked with 10-Ks. All other US mainland firms are control. Post is 1 from Q4 2024 on. Because quarterly R\&D can be noisy, I will also check yearly R\&D intensity (xrd/at) as a robustness test.
